\section*{Bab \ref{ch:b2:living-soil} --- Tanah yang Hidup}

\begin{solution}{exo:b2:living-soil:1}
O: seresah dan bahan organik yang sebagian melapuk, tempat
penguraiannya bermula. A: \omterm{def:b2:living-soil:soil}{tanah} mineral yang tercampur \omterm{def:b2:living-soil:soil}{humus}, gelap dan
berstruktur remah, yang menahan sebagian besar akar, jasad, hara, dan
airnya. B: horizon penimbunannya, tempat lempung, oksida besi, dan
karbonat yang tercuci dari atas diendapkan. C: bahan induk yang lapuk,
serpihan batuan di dalam sebuah matriks yang halus, tempat \omterm{def:b2:living-soil:soil}{tanah}
mineral yang baru dilepaskan. R: batuan dasar yang belum lapuk.
\end{solution}

\begin{solution}{exo:b2:living-soil:2}
Muatan negatif pada permukaan lempung dan \omterm{def:b2:living-soil:soil}{humusnya} menahan kation yang
dapat ditukar yang dapat diserap akarnya sebagai tukaran proton. Butir
pasirnya adalah kuarsa, yang tak bermuatan dan kasar, dengan sedikit
permukaan tiap gram dan nyaris tanpa \omterm{def:b2:living-soil:soil}{humus}. Pengapuran menambahkan
kalsium karbonat: kalsiumnya menggeser ion hidrogen dan aluminium yang
menjenuhi tapak tukar sebuah \omterm{def:b2:living-soil:soil}{tanah} yang masam, karbonatnya menetralkan
keasamannya, dan tapaknya diisi ulang dengan sebuah kation hara.
\end{solution}

\begin{solution}{exo:b2:living-soil:3}
Penguraianya: bakteri (\emph{Bacillus}) dan jamur
(\emph{Trichoderma}); perumput mikrobanya: \omterm{def:b2:unicellular-diversity:protist}{protista} (amoeba) dan
nematoda pemakan bakteri; detritivornya: \omterm{prop:b2:living-soil:community}{cacing tanah}, kaki seribu,
ekor pegas, kutu \omterm{def:b2:plant-meristems:wood}{kayu}; pemangsanya: tungau, nematoda pemangsa,
kelabang, kumbang \omterm{def:b2:living-soil:soil}{tanah}; serta tikus mondok atau celurut di puncaknya.
Perumputnya, yang $\mathrm{C:N}$-nya lebih tinggi daripada mangsa
bakterinya, mengeluarkan kelebihan nitrogennya sebagai amonium: ia
memineralkan nitrogen yang telah dikunci bakterinya.
\end{solution}

\begin{solution}{exo:b2:living-soil:4}
Arbuskular: jamurnya (sebuah glomeromisetes) masuk ke dalam sel korteks
akarnya lalu membentuk arbuskul; ditemukan pada kebanyakan herba dan
tanaman serta pada pohon tropika. Ektomikoriza: jamurnya (basidiomisetes
dan askomisetes, banyak di antaranya jamur payung) menyeludungi ujung
akarnya lalu tumbuh di antara selnya tanpa memasukinya; ditemukan pada
pohon iklim sedang dan boreal. Pada keduanya, tumbuhannya memberikan
gula dan jamurnya memberikan fosfat, nitrogen, air, hara mikro, dan
sedikit perlindungan dari patogennya.
\end{solution}

\begin{solution}{exo:b2:living-soil:5}
$L^{*} = I/k = 4/0.8 = \qty{5}{t/ha}$; \omterm{def:b2:biogeochemical-cycles:reservoir}{waktu tinggalnya} $1/k = 1.25$
tahun; kehilangan \qty{90}{\%} setelah $\ln 10/0.8 = 2.9$ tahun.
\end{solution}

\begin{solution}{exo:b2:living-soil:6}
Tiap \qty{100}{kg} karbon yang dimakan mikrobanya menahan \qty{40}{kg}
lalu memerlukan $40/8 = \qty{5}{kg}$ nitrogen. Jeraminya memasok
$100/80 =
\qty{1.25}{kg}$: imobilisasi sebesar \qty{3.75}{kg} tiap
\qty{100}{kg} karbon. Semangginya memasok $100/15 = \qty{6.7}{kg}$:
mineralisasi sebesar \qty{1.7}{kg}.
\end{solution}

\begin{solution}{exo:b2:living-soil:7}
\qty{20}{t/ha} \omterm{def:b2:living-soil:soil}{tanah} tiap tahun lewat cacingnya. \qty{20}{cm}
teratasnya berbobot $0.2\times 10^{4}\times 1.3 = \qty{2600}{t/ha}$:
130 tahun bagi satu kali lewat. Sepuluh ton tiap acre milik Darwin
adalah \qty{25}{t/ha}, yaitu orde yang sama --- ``tiap beberapa
tahun''-nya merujuk pada lapisan gembur permukaannya, yaitu beberapa
sentimeter, bukan seluruh \omterm{def:b2:living-soil:soil}{tanah} atasnya.
\end{solution}

\begin{solution}{exo:b2:living-soil:8}
Bakterinya: $10^{9}\times 10^{-12} = \qty{1}{mg}$ tiap gram. Hifanya:
volumenya $\pi(2\times 10^{-4})^{2}\times 10^{4} = \qty{1.26e-3}{cm^3}$,
massanya \qty{1.4}{mg} tiap gram. Tiap hektar ($2.6\times 10^{9}$ g):
\qty{2.6}{t} bakteri dan \qty{3.6}{t} jamur.
\end{solution}

\begin{solution}{exo:b2:living-soil:9}
$H^{*} = 1.5/0.02 = \qty{75}{t/ha}$ karbon. Dibajak: $1.5/0.04
= \qty{37.5}{t/ha}$; kehilangannya \qty{37.5}{t/ha}, dengan separuhnya
lenyap setelah $\ln 2/0.04 = 17$ tahun.
\end{solution}

\begin{solution}{exo:b2:living-soil:10}
Fosfat: $\sqrt{2\times 10^{-13}\times 8.64\times 10^{6}} =
\qty{1.3}{mm}$ dalam satu musim --- sebuah akar yang berjarak
\qty{1}{cm} dari akar berikutnya menjangkau seperdelapan \omterm{def:b2:living-soil:soil}{tanah} di
antara keduanya, sementara hifa yang berjarak \qty{0.1}{mm}
menjangkau semuanya. Nitrat: \qty{4}{cm}, yaitu lebih jauh daripada
jarak antarakarnya: ia datang ke akarnya dengan sendirinya (dan lewat
aliran massa bersama airnya), sehingga mikoriza menambahkan sedikit
bagi nitrat dan banyak bagi fosfat.
\end{solution}

\begin{solution}{exo:b2:living-soil:11}
Massanya $0.25\times 10^{4}\times 1.3 = \qty{3250}{t/ha}$. Kehilangan
bersihnya \qty{19}{t/ha} tiap tahun: umurnya 170 tahun. Karbon yang
hilang $20\times
0.03 = \qty{0.6}{t/ha}$ tiap tahun. Di bawah tanpa
olah \omterm{def:b2:living-soil:soil}{tanah} kehilangan bersihnya \qty{1}{t/ha}: 3250 tahun.
\end{solution}

\begin{solution}{exo:b2:living-soil:12}
Seresahnya berganti dalam setahun atau dua tahun, \omterm{def:b2:living-soil:soil}{humusnya} dalam lima
puluh tahun, horizon A mineralnya dalam beberapa ribu tahun
(pembentukannya pada \qty{0.1}{mm} tiap tahun), dan garamnya menimbun
dalam berpuluh tahun lalu teralirkan dalam berpuluh tahun jika ada
saluran airnya. Seorang petani melihat seresah dan tanamannya
menanggapi dalam satu musim dan \omterm{def:b2:living-soil:soil}{humusnya} dalam satu karier (penurunan
sepertiga memerlukan tiga puluh tahun; pemulihannya sama lamanya);
\omterm{prop:b2:living-soil:erosion}{erosi tanah} mineralnya tak terlihat dalam satu hayat dan tak dapat
dibalikkan dalam sepuluh hayat; garamnya terlihat dalam satu generasi
dan, tanpa saluran air, bersifat tetap. Peubah lambat \omterm{def:b2:living-soil:soil}{tanahnya} adalah
peubah yang tidak dicatat rekening tahunan mana pun.
\end{solution}

\begin{solution}{pb:b2:living-soil:1}
\textbf{1.} $0.25\times 10^{4}\times 1.3 = \qty{3250}{t/ha}$;
karbonnya $\qty{81}{t/ha}$.
\textbf{2.} $81/12 = \qty{6.8}{t/ha}$ nitrogen.
\textbf{3.} Bakterinya \qty{1}{mg/g}, $\qty{3.25}{t/ha}$; hifanya
$\pi(2\times 10^{-4})^{2}\times 2\times 10^{4}\times 1.1 =
\qty{2.8}{mg/g}$, \qty{9}{t/ha}: biomassa mikrobanya kira-kira
\qty{12}{t/ha}.
\textbf{4.} Kira-kira \qty{14}{t/ha} bahan hidup (massa segar), yaitu
sekitar \qty{17}{\%} karbon \omterm{def:b2:living-soil:soil}{humusnya} --- beberapa persen darinya
sebagai karbon, karena jasadnya kebanyakan air.
\textbf{5.} Karbon mikrobanya kira-kira separuh dari \qty{12}{t},
yaitu \qty{6}{t}; dengan berganti dua kali tiap tahun dan
\qty{40}{\%} tertahan, ia menghabiskan
$2\times 6/0.4 = \qty{30}{t}$ karbon lalu merespirasikan \qty{18}{t/ha}
karbon, yaitu \qty{66}{t} $\mathrm{CO_2}$ (sebuah taksiran yang kasar:
penahanan dan penggantiannya kasar).
\textbf{6.} Tiga kali produksi primer bersih tanamannya yang
\qty{5}{t}: taksirannya terlalu tinggi, karena anggapannya (seluruh
biomassanya berganti dua kali, tak satu pun tidur) bersifat murah hati;
\omterm{prop:b2:living-soil:humus}{respirasi tanah} yang sebenarnya di bawah sebuah tanaman berorde
produksi tanamannya sendiri, yaitu beberapa ton karbon tiap hektar tiap
tahun. \omterm{def:b2:living-soil:soil}{Tanahnya} bernapas sebanyak tumbuhan di atasnya.
\textbf{7.} Karbonnya $5\times 0.45 = \qty{2.25}{t}$; nitrogennya
$5\times
0.03 = \qty{150}{kg}$; $\mathrm{C:N} = 15$.
\textbf{8.} Memineralkan: mikrobanya memerlukan $2250\times 0.4/8 =
\qty{112}{kg}$ nitrogen dan sisanya menahan 150, sehingga
\qty{38}{kg/ha} dilepaskan secara bersih.
\textbf{9.} $1 - e^{-0.5} = \qty{39}{\%}$ setelah 3 bulan; $1 - e^{-2}
= \qty{86}{\%}$ setelah setahun.
\textbf{10.} $0.39\times 38 = \qty{15}{kg/ha}$.
\textbf{11.} Tanamannya memerlukan \qty{150}{kg}, sebagian besarnya
pada dua bulan pertamanya; sisanya memberikan \qty{15}{kg} dalam tiga
bulan dan \qty{38}{kg} seluruhnya --- sebuah tambahan, bukan sebuah
pasokan. Nilainya juga terletak pada nitrogen yang telah difiksasi
semangginya dan ia tahan di dalam \qty{150}{kg} sisanya sendiri, yang
sebagian besarnya masuk ke \omterm{def:b2:living-soil:soil}{humusnya} lalu kembali sepanjang
bertahun-tahun.
\textbf{12.} Jeraminya: karbonnya \qty{2250}{kg}, nitrogennya $2250/80 =
\qty{28}{kg}$; mikrobanya memerlukan \qty{112}{kg}, sehingga
\qty{84}{kg/ha} terimobilisasi dari \omterm{def:b2:living-soil:soil}{tanahnya} --- yaitu lebih dari
separuh kebutuhan tanamannya. Petaninya harus menambahkan nitrogen
bersama jeraminya, atau membajaknya masuk jauh sebelum menanam, atau
meninggalkannya di permukaan tempat ia terurai dengan lambat.
\textbf{13.} Dibajak: $1.2/0.025 = \qty{48}{t/ha}$; tanpa olah \omterm{def:b2:living-soil:soil}{tanah}:
$1.2/0.015 = \qty{80}{t/ha}$.
\textbf{14.} $H(t) = 80 - 32\,e^{-0.015t}$: perolehannya \qty{4.5}{t/ha}
setelah 10 tahun, \qty{17}{t/ha} setelah 50 tahun.
\textbf{15.} $\mathrm{d}H/\mathrm{d}t = 0.015\times 32 =
\qty{0.48}{t/ha}$ karbon pada tahun pertamanya, yaitu \qty{1.8}{t}
$\mathrm{CO_2}$.
\textbf{16.} $0.48\times 1.5\times 10^{9} = \qty{0.7}{GtC}$ pada tahun
pertamanya, yaitu \qty{7}{\%} pancaran fosilnya; pada tahun kelima
puluhnya lajunya telah turun menjadi $0.48\,e^{-0.75} = \qty{0.23}{t/ha}$,
yaitu \qty{0.34}{GtC}.
\textbf{17.} Perolehannya adalah pendekatan menuju sebuah \omterm{def:b2:biogeochemical-cycles:reservoir}{keadaan tunak}
yang baru: ia meluruh menjadi nol ketika $H$ mencapai \qty{80}{t/ha},
dan seluruh perolehannya dikembalikan kepada udaranya dalam berpuluh
tahun jika ladangnya dibajak lagi --- sebuah lubuk \omterm{def:b2:living-soil:soil}{tanah} adalah sebuah
pemindahan sekali jalan yang dapat balik, bukan sebuah penyeimbang yang
tetap bagi pancaran yang berlanjut.
\textbf{18.} $k_h = 0.02$: $H^{*} = \qty{60}{t/ha}$, yaitu sebuah
kehilangan \qty{20}{t/ha}, lebih banyak daripada perolehan tanpa olah
\omterm{def:b2:living-soil:soil}{tanah} lima puluh tahun pertamanya: pemanasannya dapat membatalkan
pengelolaannya.
\textbf{19.} Dibajak: $15 - 0.5 = \qty{14.5}{t/ha}$ tiap tahun; tanpa
olah \omterm{def:b2:living-soil:soil}{tanah}: \qty{1}{t/ha}.
\textbf{20.} $3250/14.5 = 220$ tahun; $3250/1 = 3250$ tahun.
\textbf{21.} Karbonnya $15\times 0.025 = \qty{0.38}{t/ha}$; nitrogennya
$0.38/12 = \qty{31}{kg/ha}$ tiap tahun.
\textbf{22.} Peluruhan \omterm{def:b2:living-soil:soil}{humusnya} pada \omterm{def:b2:biogeochemical-cycles:reservoir}{keadaan tunak} yang dibajak adalah
$k_hH^{*} =
I = \qty{1.2}{t/ha}$ tiap tahun; erosinya menambahkan
sepertiga sebanyak itu lagi, tetapi tidak seperti peluruhannya ia tidak
diimbangi masukan --- ia menguras stoknya.
\textbf{23.} Sebagian. Karbon yang tererosi yang terkubur di dalam
sebuah endapan atau sebuah \omterm{def:b2:biogeochemical-cycles:reservoir}{tandon} boleh jadi terlindung dari
peluruhannya untuk waktu yang lama (yaitu sebuah lubuk penguburan);
tetapi agregatnya pecah dalam pengangkutannya, sebagian besar karbonnya
teroksidasi di jalan, dan kesuburan ladangnya yang hilang digantikan
pupuk yang pembuatannya memancarkan karbon. Tanda global efek karbon
erosinya masih diperdebatkan.
\textbf{24.} Seresahnya $1/k = 0.5$ tahun; \omterm{def:b2:living-soil:soil}{humusnya} $1/k_h = 40$ tahun
di bawah bajaknya dan 67 tahun tanpa olah \omterm{def:b2:living-soil:soil}{tanahnya}; \omterm{def:b2:living-soil:soil}{tanah} mineralnya
$3250/0.5 =
6500$ tahun terhadap pembentukannya, dan 220 tahun terhadap
kehilangannya di bawah bajaknya.
\textbf{25.} Stok karbonnya \qty{81}{t/ha} di dalam horizon A-nya;
umurnya 220 tahun di bawah bajaknya, 3250 tahun tanpa olah \omterm{def:b2:living-soil:soil}{tanahnya};
sisa semangginya memasok kira-kira \qty{38}{kg/ha} nitrogen mineral,
dengan 15 di antaranya pada tiga bulan pertamanya.
\end{solution}
